\apendice{Detalhes do método e dos dados}{ap:metodo}

\apendicesubtitulo{Matriz do lote}

Na notação da Seção~\nameref{sec:formulacao}, $N=\sum_{(r,v)\in E_t}x_{rv}$ conta os pares escolhidos e $K=\sum_{(r,v)\in E_t}c_{rv}x_{rv}$ soma suas distâncias. Maximiza-se $N$ e, entre as soluções com esse máximo, minimiza-se $K$. Para impedir dois veículos no mesmo pedido, dois pedidos no mesmo veículo e o uso de pares inelegíveis, exige-se:
\begin{equation}
\begin{aligned}
\sum_{v\in V_t}x_{rv}&\le1 &&\quad(r\in R_t),\\
\sum_{r\in R_t}x_{rv}&\le1 &&\quad(v\in V_t),\\
x_{rv}&\in\{0,1\} &&\quad((r,v)\in E_t),\\
x_{rv}&=0 &&\quad((r,v)\notin E_t).
\end{aligned}
\label{eq:rest}
\end{equation}

A equação~\ref{eq:obj} é resolvida como uma atribuição retangular: há uma linha por pedido pendente e colunas para os veículos e para $|R_t|$ opções fictícias de \enquote{não atribuído}. Antes de montar a matriz, a distância de cada par elegível é convertida em custo inteiro, com arredondamento para metros e um índice usado no desempate determinístico. O par elegível recebe esse custo menos um bônus $M$; a coluna fictícia custa zero e o par inelegível recebe custo proibitivo. A implementação escolhe $M$ maior que a soma possível dos custos inteiros, de modo que formar mais um par prevaleça sobre qualquer diferença de distância. Só entre soluções com a mesma cardinalidade a soma dos custos de aproximação decide. As colunas fictícias garantem ao menos tantas colunas quanto linhas.

\apendicesubtitulo{Implementação do Húngaro}

A rotina usa potenciais de linhas e colunas e insere os pedidos por caminhos aumentantes. A cada passo, atualiza a menor folga das colunas ainda não visitadas e ajusta os potenciais, mantendo os custos reduzidos não negativos; os pares da solução têm custo reduzido nulo. Essa forma primal-dual de \citeonline{kuhn1955} e \citeonline{munkres1957} segue a organização por caminhos mínimos com potenciais \cite{tomizawa1971, edmondskarp1972}, com custo assintótico $O(n^2m)$ para $n$ linhas e $m$ colunas. A rotina SciPy foi usada apenas como oráculo em testes, não como política \cite{jonker1987, crouse2016}. A verificação comparou cardinalidade e soma das distâncias com enumeração exaustiva em matrizes de até $4\times4$, com o oráculo em matrizes maiores e em 2.628 lotes reais de desenvolvimento. Empates podem produzir pares diferentes com a mesma cardinalidade e distância. A otimalidade se refere ao lote.



\apendicesubtitulo{Fontes dos dados reais}

As entidades foram extraídas do banco operacional (SQL Server) e do índice de posições de ofertas (Elasticsearch). As posições dos veículos ofertados são instantâneos a cada cerca de 2 minutos; o cadastro de veículos é uma fotografia única extraída em 08/11/2026. A telemetria de terceiros existe só em fevereiro e março e \textbf{não} é usada: nenhuma quilometragem do estudo é odométrica ou validada por ela. Os arquivos derivados são armazenados em Parquet com manifesto e SHA-256 por dia; a extração depende de consultas somente leitura e não de plataformas em nuvem.

\apendicesubtitulo{Regras de construção}

todas fixadas antes da execução congelada:
\begin{alineas}
\item \textit{Elegibilidade do dia}: entidades do dia presentes, ao menos 300 instantes de posição e ao menos 200 pedidos; os 15 dias de setembro e os 25 de outubro cumpriram (mínimos de 702 instantes e 459 pedidos), sem exclusão.
\item \textit{Disponibilidade}: a presença de um veículo no registro de ofertas não prova disponibilidade, pois esse registro inclui veículos ocupados e subcobre parte da atividade. A janela de um motorista é a jornada registrada interseccionada com a presença no registro, estendida por lacunas de até 1.200 s; lacunas maiores ficam como desconhecidas e fora da janela; compromissos externos bloqueiam o motorista. A jornada usada é a realizada, informação disponível só retrospectivamente ao avaliador, comum a todas as políticas.
\item \textit{Posição inicial}: primeira posição observada do motorista na janela. Depois disso, ocupação e posição vêm das corridas que o próprio simulador atribui; as corridas históricas não entram como compromissos.
\item \textit{Parâmetros}: fator rua/reta de 1,3012 e velocidade de aproximação de 16,079 km/h, estimados na calibração de março com as distâncias de aproximação registradas pelo aplicativo em um extrato anterior dos mesmos dados (466 aproximações); raio de despacho de 4,5 km em linha reta e paciência de 180 s, ajustados naquela calibração para reproduzir o aceite e o tempo de busca observados (janela das 4 às 11 h).
\item \textit{Duração das viagens}: a duração de cada corrida com passageiro é a observada (na falta, a mediana do dia); não depende do modelo de distância.
\item \textit{Tentativas}: novas solicitações do mesmo passageiro em pouco tempo não foram deduplicadas, e cancelamentos posteriores à atribuição não são modelados.
\item \textit{Privacidade}: identificadores pseudonimizados por HMAC-SHA256 com segredo fora do repositório; só agregados são reportados.
\end{alineas}

O número médio de motoristas por dia caiu de 106,8 no treino para 74,3 em outubro. Também diminuíram os motoristas distintos com corrida encerrada (108,4 para 70,4) e os pedidos por dia (1.207 para 941). As três séries são compatíveis com mudança da operação, mas não excluem alterações de registro ou recorte. Jornadas com vínculo ambíguo entre placa e motorista (de 36 a 47 por dia) não delimitam a janela; nesses casos, ela é construída apenas com o registro de ofertas. A política recebe pedidos já chegados e o estado produzido para o instante, enquanto o avaliador reconstrói jornadas, compromissos e presença retrospectivamente.



\apendicesubtitulo{Fator rua/reta e escolha}

O km é estimado (distância em linha reta vezes 1,3012), não medido por odômetro. Como o raio e os custos do lote usam a distância em linha reta, o fator não altera a escolha em um lote fixo: multiplicar todos os custos por uma constante positiva preserva a ordenação e o ótimo da atribuição. O fator muda, porém, o km contabilizado e o tempo de aproximação, e portanto a ocupação dos veículos e a evolução do dia; quando há limites associados à distância, como um raio medido por via, muda também quem é elegível.

\apendicesubtitulo{Geografia dos cenários sintéticos}

O instantâneo do OpenStreetMap é o da consulta à API Overpass (versão 0.7.62.11) de 08/11/2026, às 15:13 UTC, com 36.350 elementos de via, licença ODbL; a malha do IBGE usa o sistema WGS84 (EPSG:4326); ambos foram registrados com SHA-256. Pontos de parada são nós viários, um por célula H3 de resolução 8, a no máximo 20 km por via, ida e volta, do centro de Vitória. As células H3 servem só para escolher os pontos de parada e definir as zonas do gerador de demanda (e para o mapa); não participam da atribuição. O mapa no material suplementar mostra que o recorte (75 células e 368 pontos de parada) cobre a ilha de Vitória e partes de Vila Velha, Cariacica e Serra, e não toda a região metropolitana; é um recorte experimental, não um limite administrativo nem a área de atuação de uma empresa.



O arquivo do OpenStreetMap contém 521 trechos marcados como ponte ou viaduto, entre eles a Ponte Deputado Darcy Castello de Mendonça, a Ponte Governador Gerson Camata, a Ponte do Príncipe e a Ponte Florentino Avidos. As verificações feitas foram estruturais: o grafo do recorte é conexo (as distâncias entre os 368 pontos são todas finitas), a distância por via nunca é menor que a distância em linha reta, e a razão entre as duas, em mediana, é de 1,45 entre Vitória e Vila Velha e de 1,34 dentro de Vila Velha, de modo que há travessia mapeada entre a ilha e o continente. \textbf{Não} se verificou ponte a ponte se cada ligação existente está mapeada nem se validaram tempos de percurso contra medições; não se afirma, portanto, que as pontes estejam corretamente representadas, e congestionamento e travessias bloqueadas não são modelados.

\apendicesubtitulo{Gerador e variantes}

O gerador (versão 1.0.0) tem parâmetros registrados; as jornadas de 9 horas começam às 6, 8, 10 ou 12 h. Velocidade, raio e paciência são hipóteses para uma operação sintética, não calibradas. Das 17 variantes, quatro alteram a relação frota/demanda (4 ou 14 pedidos por veículo, contra 8), quatro alteram a resolução das zonas do gerador de demanda e nove alteram só parâmetros de um planejador de reposicionamento que não participa do despacho; estas nove são novos sorteios do cenário-base (o código do cenário entra na semente), e não condições diferentes. As concentrações de demanda são hipóteses e não representam padrões reais. As sementes finais (20001--20040) foram reservadas antes da execução; as de desenvolvimento (101--120) foram usadas antes.

\apendicesubtitulo{Distâncias e estatística}

Nos dados históricos, o custo de aproximação de um par é $c_{rv}=1{,}3012\,h_{rv}$, em que $h_{rv}$ é a distância em linha reta. Nos cenários sintéticos, $c_{rv}$ é a distância pelo caminho mínimo na rede viária; o fator histórico não se aplica. A diferença em km é $\Delta K=K_C-K_{ref}$. A variação relativa de atendidos, $100\,\Delta N/N_{ref}$, usa os atendidos da referência, ao passo que a diferença de taxas em pontos percentuais usa todos os pedidos elegíveis. As medidas não são intercambiáveis.

A espera inclui o intervalo entre chegada e despacho e o deslocamento estimado até a origem, entre os pedidos atendidos de cada política. O IC95\% da diferença média diária usa \textit{bootstrap} de blocos móveis circulares de 7 dias, com 10.000 reamostragens e semente 20261108; intervalos alternativos constam do material suplementar. Um intervalo que inclui zero não demonstra equivalência ou não inferioridade, pois nenhuma margem foi fixada previamente. Nos cenários sintéticos, os 20 testes previstos no plano (5 cenários principais $\times$ 4 variantes de despacho em lote) usam correção de Holm. Como seus resultados já eram conhecidos na decisão de adotar o método, sua interpretação no artigo é exploratória.


\apendice{Verificações, análises posteriores e reprodutibilidade}{ap:repro}

O congelamento local anterior à simulação de setembro e outubro registrou o SHA-256 de 151 arquivos do código e dos \textit{scripts} de avaliação (\texttt{480a98a1\ldots e93f87}), o \textit{commit} de base \texttt{3e77cb9}, Python 3.12.14, a configuração das três políticas e 201 manifestos de dados com SHA-256 conferidos contra os Parquet. Naquele estado do código, os testes automáticos tiveram 535 aprovações e 1 salto; a análise estática não apontou erros. Esses registros locais permitem conferir arquivos, mas não certificam externamente a data do congelamento.

Em quatro dias de desenvolvimento (10/02, 04/03, 17/03 e 25/03), o Húngaro coincidiu com o oráculo em cardinalidade e soma de distâncias, à resolução de 1 m, nos 2.628 lotes examinados. Também foram conferidos chegada anterior à decisão, instantes múltiplos de 60 s, pares únicos e elegíveis e ausência de serviços simultâneos por motorista. O lote sequencial coincidiu com uma reimplementação independente. Todos os 15 dias de setembro e os 25 de outubro foram elegíveis. Houve 19 e 17 compromissos, respectivamente, com posição de entrada inconsistente; essas inconsistências eram iguais nas três políticas e não foram atribuídas a elas.

Depois da avaliação, o código foi reorganizado, alterando seu hash sem reescrever os manifestos congelados. Reexecuções amostrais de três pares sintéticos e um dia histórico reproduziram os resultados anteriores; a análise dos arquivos congelados reproduziu as tabelas. Em seguida, as três políticas reexecutadas nas entradas congeladas reproduziram os 120 resultados de dia e política de setembro e outubro, e a reconstrução atual das entradas reproduziu os totais dos dois períodos. Essas verificações de equivalência não são novos testes científicos. Na rede viária sintética, foram conferidas conectividade e relações entre distâncias por via e em linha reta; não houve validação dos tempos de percurso nem inspeção de cada ponte.







As tabelas detalhadas de validação operacional e recalibração, juntamente com as séries diárias, os intervalos por variante, os cortes mensais, a conferência da coleta, os resultados completos dos cenários sintéticos e os caminhos dos arquivos, estão no \caminho{article/supplement/material_suplementar.pdf}.
